%% 自動生成（tools/mkchecklist.py）。各章の「ふりかえり」の項目を集めたもの
\begingroup\footnotesize
\def\bkcb{\tikz[baseline=-0.5ex]\draw[\bkTF, line width=0.7pt, rounded corners=0.5pt] (0,-0.3em) rectangle (0.8em,0.5em);}
\begin{longtblr}{colspec={Q[c,m,wd=9mm] X[l,m] Q[c,m,wd=10mm] Q[c,m,wd=9mm] Q[c,m,wd=9mm]}, width=\linewidth, rowhead=1, colsep=3pt,
  row{1}={bg=bkVol!14, font=\sffamily\bfseries\scriptsize}, hlines={0.4pt,bkLine}, rowsep=0.7pt}
章 & できるようになったこと & ページ & 1回め & 2回め \\
{\sffamily\bfseries\color{\bkTT}7-1} & 数の表を行列として読み、成分 $a_{ij}$ を言える。和・差・実数倍を計算できる。 & \pageref{bk:7-1:gyoretsu} & \bkcb & \bkcb \\
 & 行列×ベクトル・行列×行列を、行と列をなぞって計算できる。 & \pageref{bk:7-1:kakeru} & \bkcb & \bkcb \\
 & 連立方程式を $A\vec x=\vec b$ の形に書き、表計算や Python で確かめられる。 & \pageref{bk:7-1:renritsu} & \bkcb & \bkcb \\
{\sffamily\bfseries\color{\bkTT}7-2} & 行列×ベクトルを、行と列をなぞって計算できる。 & \pageref{bk:7-2:nazoru} & \bkcb & \bkcb \\
 & 列（物差しの矢印の行き先）から、絵の変わり方を読みとり、行き先から行列を作れる。 & \pageref{bk:7-2:retsu} & \bkcb & \bkcb \\
 & ヤコビ行列や応力の行列を、「入れると出てくる機械」として計算できる。 & \pageref{bk:7-2:jacobi} & \bkcb & \bkcb \\
{\sffamily\bfseries\color{\bkTT}7-3} & 行列×行列を、行と列の内積で計算できる。 & \pageref{bk:7-3:naiseki} & \bkcb & \bkcb \\
 & 変換をつなぐと積になり、順番で結果が変わることを説明できる。 & \pageref{bk:7-3:junban} & \bkcb & \bkcb \\
 & 電気の箱をつないだ回路を、行列の積で計算できる。 & \pageref{bk:7-3:hako} & \bkcb & \bkcb \\
{\sffamily\bfseries\color{\bkTT}7-4} & 方眼で面積を数え、$ad-bc$ で行列式を計算できる。 & \pageref{bk:7-4:adbc} & \bkcb & \bkcb \\
 & 行列式のマイナス（裏返し）と0（平らにつぶれる）の意味がわかる。 & \pageref{bk:7-4:zero} & \bkcb & \bkcb \\
 & 3×3 の行列式をサラスの方法で計算し、工学の行列式を読める。 & \pageref{bk:7-4:sanjigen} & \bkcb & \bkcb \\
{\sffamily\bfseries\color{\bkTT}7-5} & 2×2 の逆行列を作って確かめ、$\vec x=A^{-1}\vec b$ で連立方程式を解ける。 & \pageref{bk:7-5:gyaku} & \bkcb & \bkcb \\
 & 掃き出し法で、3つの文字の連立方程式と逆行列を求められる。 & \pageref{bk:7-5:hakidashi} & \bkcb & \bkcb \\
 & 正規方程式 $A^{\mathsf T}A\vec x=A^{\mathsf T}\vec b$ で、最小二乗法の直線を求められる。 & \pageref{bk:7-5:saisho} & \bkcb & \bkcb \\
{\sffamily\bfseries\color{\bkTT}7-6} & 回転行列 $R(\theta)$ を作り、点や図形を回せる。 & \pageref{bk:7-6:tsukuru} & \bkcb & \bkcb \\
 & $R(\alpha)R(\beta)=R(\alpha+\beta)$ と $R(\theta)^{-1}=R(\theta)^{\mathsf T}$ を使える。 & \pageref{bk:7-6:tsunagu} & \bkcb & \bkcb \\
 & 軸を回して点や力を見たり、腕の手の位置を求めたりできる。 & \pageref{bk:7-6:jiku} & \bkcb & \bkcb \\
{\sffamily\bfseries\color{\bkTT}7-7} & 3本の軸のまわりの回転を $R_x$・$R_y$・$R_z$ で表し、点を回せる。 & \pageref{bk:7-7:rz} & \bkcb & \bkcb \\
 & ロール・ピッチ・ヨーから機首の向きを求め、順番で向きが変わることを説明できる。 & \pageref{bk:7-7:euler} & \bkcb & \bkcb \\
 & 同次変換行列で、腕の手の位置やカメラで見た物の位置を求められる。 & \pageref{bk:7-7:doji} & \bkcb & \bkcb \\
{\sffamily\bfseries\color{\bkTT}7-8} & 行列をかけても向きが変わらない矢印を、図と表計算でさがせる。 & \pageref{bk:7-8:sagasu} & \bkcb & \bkcb \\
 & $A\vec v$ を計算して固有ベクトルを確かめ、振動モードが固有ベクトルだと説明できる。 & \pageref{bk:7-8:daisha} & \bkcb & \bkcb \\
 & 特性方程式で固有値と固有ベクトルを求め、6-12 の $s$ と同じものだと確かめられる。 & \pageref{bk:7-8:tokusei} & \bkcb & \bkcb \\
{\sffamily\bfseries\color{\bkTT}7-9} & 固有ベクトルの物差しで見ると「のびるだけ」になることを、図で説明できる。 & \pageref{bk:7-9:naname} & \bkcb & \bkcb \\
 & $P$ と $D$ を作り、$P^{-1}AP=D$ を計算で確かめられる。 & \pageref{bk:7-9:jiku} & \bkcb & \bkcb \\
 & 対角化で $A^n$ を計算したり、台車の式をモードの座標でほどいたりできる。 & \pageref{bk:7-9:hodoku} & \bkcb & \bkcb \\
{\sffamily\bfseries\color{\bkTT}7-10} & 同じ行列を何回もかけて、点の向かう先を表計算と図で調べられる。 & \pageref{bk:7-10:kurikaeshi} & \bkcb & \bkcb \\
 & 倍々のくり返しと、等比級数の合計を計算できる。 & \pageref{bk:7-10:kyusu} & \bkcb & \bkcb \\
 & 固有値の大きさで行き先を見分け、ばねがふくらむわけを説明できる。 & \pageref{bk:7-10:bane} & \bkcb & \bkcb \\
{\sffamily\bfseries\color{\bkTT}7-11} & ロボットの腕を、ヤコビ行列の逆行列で目標へ近づけるしくみを説明できる。 & \pageref{bk:7-11:gyaku} & \bkcb & \bkcb \\
 & ばねの網と抵抗の網、箱をつないだ回路を、行列で書ける。 & \pageref{bk:7-11:ami} & \bkcb & \bkcb \\
 & 長い電線が $\cosh$・$\sinh$ の行列になるわけを、箱のくさりとして読める。 & \pageref{bk:7-11:bunpu} & \bkcb & \bkcb \\
\end{longtblr}
\endgroup
